\section{Performance Evaluation}

\label{sec:evaluation}

\subsection{Signing Latency}

\begin{table}[t]
\centering
\begin{tabular}{lcccc}
\toprule
\textbf{Protocol} & \textbf{Threshold} & \textbf{Keygen (ms)} & \textbf{Pre-sign (ms)} & \textbf{Sign (ms)} \\
\midrule
CGGMP21 & (2, 3) & 850 & 120 & 45 \\
CGGMP21 & (3, 5) & 2,100 & 340 & 95 \\
CGGMP21 & (5, 9) & 8,400 & 1,200 & 280 \\
FROST & (2, 3) & 180 & --- & 12 \\
FROST & (3, 5) & 420 & --- & 28 \\
FROST & (5, 9) & 1,600 & --- & 85 \\
\bottomrule
\end{tabular}
\caption{MPC protocol latency on Zoo testnet. FROST is 3--4$\times$ faster than CGGMP21 for signing.}
\label{tab:latency}
\end{table}

\subsection{Throughput}

\begin{table}[t]
\centering
\begin{tabular}{lcc}
\toprule
\textbf{Protocol} & \textbf{Signs/sec (2,3)} & \textbf{Signs/sec (3,5)} \\
\midrule
CGGMP21 & 22 & 10 \\
CGGMP21 (pre-signed) & 420 & 180 \\
FROST & 83 & 35 \\
\bottomrule
\end{tabular}
\caption{Signing throughput. CGGMP21 with pre-signing achieves highest throughput for batch operations.}
\label{tab:throughput}
\end{table}

\subsection{Network Overhead}

\begin{table}[t]
\centering
\begin{tabular}{lccc}
\toprule
\textbf{Phase} & \textbf{Messages} & \textbf{Bytes/Party} & \textbf{Rounds} \\
\midrule
CGGMP21 Keygen & $O(n^2)$ & 12,800 & 5 \\
CGGMP21 Pre-sign & $O(t^2)$ & 4,200 & 4 \\
CGGMP21 Sign & $O(t)$ & 256 & 1 \\
FROST Keygen & $O(n^2)$ & 3,200 & 2 \\
FROST Sign & $O(t)$ & 128 & 2 \\
\bottomrule
\end{tabular}
\caption{Network communication overhead for (2, 3)-threshold.}
\label{tab:network}
\end{table}

\subsection{Security Analysis}

\begin{theorem}[Key Security]
Under the $(2, 3)$-threshold, an adversary controlling a single party (researcher, DAO, or escrow) cannot produce a valid signature or learn any information about the private key.
\end{theorem}

\begin{proof}
Both CGGMP21 (UC-secure~\cite{cggmp2021}) and FROST (secure in the algebraic group model~\cite{komlo2020frost}) guarantee that $t - 1$ shares reveal no information about $sk$ under the discrete logarithm assumption. With $t = 2$ and a single corrupted party, the adversary holds 1 share, which is insufficient.
\end{proof}

\begin{theorem}[Identifiable Abort]
If a signing session fails due to a malicious party, the protocol identifies the malicious party with probability 1.
\end{theorem}

CGGMP21 provides identifiable abort natively. FROST achieves it via the identifiable abort extension, where each party proves correctness of their partial signature via a Schnorr proof of knowledge.

